% Part: lambda-calculus
% Chapter: introduction
% Section: fixed-point-combinator

\documentclass[../../../include/open-logic-section]{subfiles}

\begin{document}

\olfileid{lam}{int}{fix}
\olsection{স্থির-বিন্দু সমাবেশক}

ধরা যাক, তোমার কাছে একটি ল্যাম্বডা পদ~$g$ আছে, এবং এমন আরেকটি পদ~$k$
চাই যার ধর্ম হলো, $k$ পদটি $gk$-এর $\beta$-সমতুল্য। আমাদের সংকেতলিপির
প্রথা অনুসারে পদদুটি সংজ্ঞায়িত করো:
\[
\fn{diag}(x) = xx
\]
এবং
\[
l(x) = g(\fn{diag}(x))
\]
অর্থাৎ $l$ হলো $\lambd[x][g(xx)]$ পদ। $k$-কে $ll$ পদটি ধরো। তাহলে
\begin{align*}
k & = (\lambd[x][g(xx)])(\lambd[x][g(xx)]) \\
& \red  g((\lambd[x][g(xx)])(\lambd[x][g(xx)])) \\
& = gk.
\end{align*}
যদি নেওয়া হয়
\[
Y = \lambd[g][((\lambd[x][g(xx)])(\lambd[x][g(xx)]))]
\]
তাহলে $Yg$ ও $g(Yg)$ একটি সাধারণ পদে হ্রাস পায়; তাই $Yg \equiv_\beta
g(Yg)$। একে ``কারির সমাবেশক'' বলা হয়। এর বদলে যদি নেওয়া হয়
\[
Y = (\lambd[xg][g(xxg)])(\lambd[xg][g(xxg)])
\]
তবে বস্তুত $Yg$, $g(Yg)$-তে হ্রাস পায়—এটি আরও শক্তিশালী বক্তব্য।
$Y$-এর এই শেষের রূপটিকে ``টুরিংয়ের সমাবেশক'' বলা হয়।

\end{document}
